\chapter{The Routing Table and How a Router Decides}\label{ch:routingtable}
\bp{3.1}

\begin{outcomes}
By the end of this chapter you will be able to:
\begin{tight}
  \item \textbf{Read} every column and code in \cmd{show ip route}.
  \item \textbf{Apply} the longest prefix match rule to determine which entry a
        given packet uses.
  \item \textbf{Explain} administrative distance and predict which source wins
        when two protocols offer the same prefix.
  \item \textbf{Distinguish} the three reasons a packet is dropped: no route, a
        route to null, and an unreachable next hop.
\end{tight}
\end{outcomes}

\begin{prereq}
\chref{ipv4} for prefixes and masks --- this chapter is unreadable without
fluent subnetting. \chref{trunking} for SVIs as gateways.
\end{prereq}

\begin{keyterms}
routing table $\bullet$ prefix $\bullet$ next hop $\bullet$ exit interface
$\bullet$ administrative distance $\bullet$ metric $\bullet$ longest prefix match
$\bullet$ directly connected $\bullet$ local route $\bullet$ default route
$\bullet$ gateway of last resort $\bullet$ null0
\end{keyterms}

\section{Reading the table}

\begin{verify}{R1}{show ip route}
Codes: L - local, C - connected, S - static, R - RIP, M - mobile, B - BGP
       O - OSPF, IA - OSPF inter area, N1 - OSPF NSSA external type 1
       E1 - OSPF external type 1, E2 - OSPF external type 2
       * - candidate default, U - per-user static route

Gateway of last resort is 203.0.113.1 to network 0.0.0.0

S*    0.0.0.0/0 [1/0] via 203.0.113.1
      10.0.0.0/8 is variably subnetted, 4 subnets, 2 masks
O        10.1.0.0/16 [110/20] via 10.255.0.2, 00:14:22, GigabitEthernet0/1
C        10.255.0.0/30 is directly connected, GigabitEthernet0/1
L        10.255.0.1/32 is directly connected, GigabitEthernet0/1
S        10.9.9.0/24 [1/0] via 10.255.0.2
      192.168.10.0/24 is directly connected, Vlan10
C        192.168.10.0/24 is directly connected, Vlan10
L        192.168.10.1/32 is directly connected, Vlan10
\end{verify}

\begin{longtable}{@{}L{3.0cm}L{12.2cm}@{}}
\toprule
\rowcolor{primary!15}
\textbf{Element} & \textbf{What it tells you} \\
\midrule
\endhead
Code letter & Where the route came from. \cmd{C} connected, \cmd{L} local,
\cmd{S} static, \cmd{O} OSPF \\
\rowcolor{lightbg}
Prefix and mask & The destination network this entry covers \\
\cmd{[110/20]} & \textbf{Administrative distance / metric.} 110 is OSPF's
distance; 20 is OSPF's cost for this path \\
\rowcolor{lightbg}
\cmd{via 10.255.0.2} & The next hop --- the neighbouring router to hand the
packet to \\
Interface & The exit interface \\
\rowcolor{lightbg}
\cmd{00:14:22} & How long the route has been in the table. A route that keeps
resetting is flapping \\
\bottomrule
\end{longtable}

\begin{conceptbox}[C and L are not the same thing]
\cmd{C} is the \emph{subnet} on a directly connected interface --- the range the
router can reach by ARP. \cmd{L} is a \cmd{/32} host route for the router's
\emph{own} address on that interface, so that traffic addressed to the router
itself is delivered locally rather than forwarded.

You do not configure \cmd{L} routes; the router creates one for every address you
give an interface. They appear in every table and are not a fault.
\end{conceptbox}

\section{Rule one: longest prefix match}

\begin{conceptbox}[The most specific route wins, regardless of anything else]
When several entries could carry a packet, the router uses the one with the
\textbf{longest prefix} --- the most 1 bits in the mask. Administrative distance
and metric never override this; they only choose between entries with the
\emph{same} prefix length.
\end{conceptbox}

\begin{worked}[Which entry carries the packet?]
Using the table above, a packet for \cmd{10.9.9.20}:

\begin{tightnum}
  \item \cmd{0.0.0.0/0} --- matches everything. Prefix length 0.
  \item \cmd{10.1.0.0/16} --- does \cmd{10.9.9.20} fall inside? The /16 covers
        \cmd{10.1.0.0}--\cmd{10.1.255.255}. No.
  \item \cmd{10.9.9.0/24} --- covers \cmd{10.9.9.0}--\cmd{10.9.9.255}. Yes.
        Prefix length 24.
\end{tightnum}

Two entries match: the default and the /24. The /24 is longer, so the packet goes
\cmd{via 10.255.0.2}. Note that the default route has a \emph{better}
administrative distance (1 versus 1 --- equal here) and it still loses, because
prefix length is decided first.

Now a packet for \cmd{172.20.5.1}: only \cmd{0.0.0.0/0} matches, so it takes the
gateway of last resort.
\end{worked}

\section{Rule two: administrative distance}

\begin{definitionbox}[Administrative distance]
\term{Administrative distance} ranks the trustworthiness of a route's
\emph{source}. When two sources offer the same prefix at the same length, the
lower distance is installed and the other is discarded --- it does not even
appear in the table.
\end{definitionbox}

\begin{longtable}{@{}L{5.4cm}C{2.6cm}L{7.2cm}@{}}
\toprule
\rowcolor{primary!15}
\textbf{Source} & \textbf{Distance} & \textbf{Note} \\
\midrule
\endhead
Connected & 0 & Nothing beats seeing it yourself \\
\rowcolor{lightbg}
Static & 1 & You said so, so the router believes you \\
eBGP & 20 & \\
\rowcolor{lightbg}
OSPF & 110 & The only dynamic protocol in this exam \\
RIP & 120 & Legacy \\
\rowcolor{lightbg}
Unreachable / unusable & 255 & Never installed \\
\bottomrule
\end{longtable}

\begin{conceptbox}[Why a static route beats OSPF, and when that bites]
A static route at distance 1 will always displace an OSPF route at 110 for the
same prefix. That is usually what you want --- until the static route's path
fails. OSPF would have re-routed; the static route sits there pointing at a dead
path, because a static route has no way of knowing anything is wrong beyond the
directly connected interface.

The deliberate use of this is the \term{floating static route}: a backup static
configured with a distance \emph{higher} than the dynamic protocol, so it stays
out of the table until the dynamic route disappears. That is \chref{static}.
\end{conceptbox}

\section{Three ways a packet dies}

\begin{longtable}{@{}L{3.6cm}L{5.2cm}L{6.2cm}@{}}
\toprule
\rowcolor{primary!15}
\textbf{Cause} & \textbf{What the router does} & \textbf{How you see it} \\
\midrule
\endhead
No matching route & Drops the packet and sends ICMP unreachable &
\cmd{U} in ping output; \cmd{show ip route} has no entry and no default \\
\rowcolor{lightbg}
Route points to \cmd{Null0} & Silently discards & An entry
\cmd{via Null0}. Often deliberate, for summarisation or blackholing \\
Next hop unreachable & Route exists but cannot be resolved to a next hop &
Route is in the table; pings fail; \cmd{show ip route} for the next hop shows
nothing \\
\bottomrule
\end{longtable}

\begin{alertbox}[A route in the table is not proof of reachability]
The presence of an entry means the router knows where it \emph{would} send the
packet. It says nothing about whether the next hop is alive, whether the far end
has a route back, or whether an access list drops it en route.

Confirm with an extended ping sourced from the relevant interface
(\chref{transport}) --- that tests the return path too, which a route entry never
does.
\end{alertbox}

\begin{verify}{R1}{show ip route 10.9.9.20}
Routing entry for 10.9.9.0/24
  Known via "static", distance 1, metric 0
  Routing Descriptor Blocks:
  * 10.255.0.2
      Route metric is 0, traffic share count is 1
\end{verify}

Querying a specific address rather than reading the whole table is faster and is
what the exam expects when it asks which entry would be used.

\begin{pitfall}
\begin{tight}
  \item \textbf{Comparing metrics across protocols.} An OSPF cost of 20 and a RIP
        hop count of 2 are not comparable. Administrative distance decides first,
        and only then does metric compare \emph{within} one protocol.
  \item \textbf{Thinking a better distance beats a longer prefix.} It does not.
        Prefix length is decided first, always.
  \item \textbf{Missing the gateway of last resort line.} If it says
        \cmd{not set}, there is no default and anything unmatched is dropped.
  \item \textbf{Reading \cmd{L} entries as duplicates.} They are the router's own
        addresses and belong there.
  \item \textbf{Assuming symmetry.} Traffic reaching a destination does not mean
        the reply can return. Most one-way faults are a missing reverse route.
\end{tight}
\end{pitfall}

\begin{breakfix}[The route is there and the ping still fails.]
Users on \cmd{192.168.10.0/24} cannot reach \cmd{10.9.9.0/24}. \cmd{R1} has the
route:

\begin{verify}{R1}{show ip route 10.9.9.0}
Routing entry for 10.9.9.0/24
  Known via "static", distance 1, metric 0
  Routing Descriptor Blocks:
  * 10.255.0.2
\end{verify}

A ping from \cmd{R1} to \cmd{10.9.9.20} succeeds. A ping from a user's PC fails.

\textbf{Diagnosis.} The difference between the two pings is the source address.
\cmd{R1}'s own ping is sourced from \cmd{10.255.0.1}, which the far end knows
about because it is the link between them. The PC's traffic is sourced from
\cmd{192.168.10.x}, and the far router has no route back to that subnet. The
forward path is fine; the return path does not exist.

\textbf{Confirm it} without leaving \cmd{R1}, using an extended ping that borrows
the user's source address:

\begin{verify}{R1}{ping 10.9.9.20 source Vlan10}
Sending 5, 100-byte ICMP Echos to 10.9.9.20, timeout is 2 seconds:
Packet sent with a source address of 192.168.10.1
.....
Success rate is 0 percent (0/5)
\end{verify}

Same destination, different source, opposite result --- which is proof, not
inference.

\textbf{Fix.} Add a route on the far router back to \cmd{192.168.10.0/24}. Then
re-run the same extended ping and expect five exclamation marks.
\end{breakfix}

\begin{lab}{Read the table, then predict the path}{Packet Tracer}
\textbf{Topology.} Three routers in a line, a LAN on each end, and a second path
between the outer two.

\begin{tasks}
  \item Address everything and confirm the connected routes appear. Identify
        every \cmd{C} and \cmd{L} entry and explain each \cmd{L}.
  \item Add static routes so the two LANs can reach each other. Record the
        distance and metric shown.
  \item Add a default route on one router. Note the \cmd{S*} and the
        \emph{gateway of last resort} line.
  \item Write down, before testing, which entry a packet to each of four
        addresses would use. Verify each with \cmd{show ip route <address>}.
  \item Add a second, more specific route for one destination and confirm it wins
        over the less specific one, regardless of distance.
  \item Remove the reverse route on the far router and reproduce the one-way
        fault above using \cmd{ping ... source}.
  \item \textbf{Extension.} Add a route pointing to \cmd{Null0} for an unused
        prefix and observe how traffic to it fails compared with having no route
        at all.
\end{tasks}
\end{lab}

\begin{checkpoint}
\begin{tightnum}
  \item A table contains \cmd{10.0.0.0/8} via OSPF and \cmd{10.5.0.0/16} via a
        static route. Which carries a packet to \cmd{10.5.1.1}, and why?
  \item Both OSPF and a static route offer \cmd{192.168.50.0/24}. Which is
        installed, and what happens to the other?
  \item A route exists but pings fail. Name three causes and the single command
        that distinguishes the most common one.
\end{tightnum}
\end{checkpoint}

\begin{chaptersummary}
\begin{tight}
  \item One table, many sources. \cmd{C} connected, \cmd{L} the router's own
        address, \cmd{S} static, \cmd{O} OSPF.
  \item \cmd{[distance/metric]}. Distance ranks the source; metric ranks paths
        within one protocol and is never compared across protocols.
  \item Longest prefix match decides first. Distance only breaks ties between
        equal-length prefixes.
  \item Connected 0, static 1, OSPF 110. A static route displaces OSPF and will
        not fail over unless you make it float.
  \item A route in the table proves intent, not reachability. Test the return
        path with an extended ping.
\end{tight}
\end{chaptersummary}

\begin{reviewq}
  \item \textbf{(MCQ)} A router has \cmd{0.0.0.0/0}, \cmd{172.16.0.0/16} and
        \cmd{172.16.8.0/24}. Which carries a packet to \cmd{172.16.8.50}?
        (a)~the default (b)~the /16 (c)~the /24 (d)~it is dropped
  \item \textbf{(MCQ)} What is the administrative distance of a static route?
        (a)~0 (b)~1 (c)~110 (d)~120
  \item \textbf{(Short)} Explain the difference between a \cmd{C} and an \cmd{L}
        entry and why both exist.
  \item \textbf{(Short)} Why can an OSPF cost never be compared with a RIP hop
        count?
  \item \textbf{(Applied)} A router shows a static route to a remote subnet, and
        pings from the router succeed while pings from its LAN fail. Explain the
        likely cause, the exact command that would prove it, and the fix ---
        including which device you would change.
\end{reviewq}
